\chapter{Toroidal Compactification and the Narain Lattice}\label{ch:narain}

A string on a circle carries two integers, the momentum number $n$ and the winding number
$w$, and \cref{ch:circle,ch:tduality} showed that the whole radius dependence of the spectrum
sits in the way these two integers are split into a left-moving and a right-moving momentum.
This chapter asks what happens when $d$ dimensions are curled up at once, on a torus $T^d$
with an arbitrary constant metric $G_{ij}$ and two-form $B_{ij}$. The $2d$ integers of
momentum and winding form a lattice $\Gamma^{d,d}$ with an inner product of signature $(d,d)$;
it is \emph{even} and \emph{self-dual}, and as an abstract lattice it is the same for every
torus. All the dependence on the $d^2$ parameters $(G,B)$ is in the choice of a
$d$-dimensional positive subspace of $\R^{d,d}$, and the set of such choices is the
homogeneous space $\OddR{d}/\big(O(d)\times O(d)\big)$. A student should care for three
reasons. First, the duality group $\Odd{d}$ of \cref{ch:odd}, the generalized metric of
\cref{ch:genmetric} and the doubled coordinates of \cref{ch:doubled} are all read off from
this lattice. Second, the same two words, even and self-dual, describe the cohomology lattice
of a K3 surface (\cref{ch:lattices,ch:k3lattice}) and organise the moduli space of strings on
K3 (\cref{ch:stringsk3}). Third, the rank-two case is small enough to be checked by the Lean
kernel; \cref{sec:narain:lean} states exactly what was checked.

\section{The torus and its background fields}\label{sec:narain:torus}

\subsection{Coordinates of period \texorpdfstring{$2\pi$}{2 pi}}

A flat torus $T^d$ is the quotient of $\R^d$ by a lattice. Let $\hat e_1,\dots,\hat e_d$ be
$d$ linearly independent vectors in $\R^d$, with the dimension of a length, and identify
\begin{equation}\label{eq:narain:torusdef}
 x\;\sim\;x+2\pi\sum_{i=1}^{d}m^i\,\hat e_i,\qquad m^i\in\Z .
\end{equation}
For $d=1$ and $\hat e_1=R$ this is the circle $X\sim X+2\pi R$ of \cref{ch:circle}. Instead of
Cartesian coordinates it is convenient to use coordinates \emph{along the lattice vectors}:
write $x=\sum_iX^i\hat e_i$. The numbers $X^i$ are dimensionless, and the identification
\eqref{eq:narain:torusdef} becomes
\begin{equation}\label{eq:narain:period}
 X^i\;\sim\;X^i+2\pi m^i ,\qquad m^i\in\Z,
\end{equation}
the same for every torus \source{GPR §2.4, eq.~(2.4.2), ll.~1055--1058}. The price is that
the metric is no longer the unit matrix. The squared length of a displacement is
$\dd s^2=\hat G_{ij}\,\dd X^i\dd X^j$ with $\hat G_{ij}=\hat e_i\cdot\hat e_j$, a constant,
symmetric, positive definite $d\times d$ matrix with the dimension of length squared. Shape
and size of the torus are encoded in $\hat G_{ij}$, and the coordinate ranges are fixed once
and for all (\cref{fig:narain:torus}).

Because T-duality involves $\ap$ explicitly, it pays to measure lengths in string units. We
define the dimensionless metric
\begin{equation}\label{eq:narain:Gdimless}
 G_{ij}=\frac{\hat G_{ij}}{\ap}=\frac{\hat e_i\cdot\hat e_j}{\ap}
\end{equation}
\source{GPR §2.2, ll.~805--821}. For the circle, $G_{11}=R^2/\ap$; the self-dual radius
$R=\sqrt{\ap}$ of \cref{ch:tduality} is $G_{11}=1$. Until \cref{sec:narain:massformula}
restores $\ap$ in the mass formula, every quantity in this chapter is dimensionless.

\begin{figure}[t]
\centering
\begin{tikzpicture}[scale=0.78,>=Stealth]
 % lattice points
 \foreach \a in {0,1,2,3}{\foreach \b in {0,1,2}{
   \fill[black!55] ({2.2*\a+0.8*\b},{1.5*\b}) circle (1.3pt);}}
 % fundamental cell
 \fill[blue!8] (0,0)--(2.2,0)--(3.0,1.5)--(0.8,1.5)--cycle;
 \draw[black!40] (0,0)--(6.6,0) (0.8,1.5)--(7.4,1.5) (1.6,3.0)--(8.2,3.0);
 \foreach \a in {0,1,2,3}{\draw[black!40] ({2.2*\a},0)--({2.2*\a+1.6},3.0);}
 \draw[->,very thick,blue!60!black] (0,0)--(2.2,0) node[midway,below]{$2\pi\hat e_1$};
 \draw[->,very thick,blue!60!black] (0,0)--(0.8,1.5) node[midway,left]{$2\pi\hat e_2$};
 % a winding string with m=(1,1)
 \draw[red!70!black,thick,decorate,decoration={snake,amplitude=0.6mm,segment length=5mm}]
   (2.6,0.35)--(5.6,1.85);
 \fill[red!70!black] (2.6,0.35) circle (1.6pt) (5.6,1.85) circle (1.6pt);
 \node[red!70!black,fill=white,inner sep=1pt] at (6.6,2.25) {$m=(1,1)$};
 \node[black!70] at (1.55,0.75) {$T^2$};
\end{tikzpicture}
\caption{A two-torus is the plane modulo the lattice generated by $2\pi\hat e_1$ and
$2\pi\hat e_2$ (shaded: one cell). In the coordinates $X^i$ the cell is always
$[0,2\pi)^2$, and the geometry sits in $G_{ij}=\hat e_i\cdot\hat e_j/\ap$. Wavy line: a closed
string with winding numbers $m=(1,1)$.}
\label{fig:narain:torus}
\end{figure}

\subsection{The action and its \texorpdfstring{$d^2$}{d squared} parameters}

\Cref{ch:backgrounds} showed that a string couples to a background metric and to a background
two-form, the Kalb--Ramond field\index{Kalb--Ramond field}\index{B-field}. On a torus both may
be taken constant. With a flat worldsheet metric, worldsheet coordinates $(\tau,\sigma)$,
$\sigma\in[0,2\pi)$, and the notation $\dot X=\partial_\tau X$, $X'=\partial_\sigma X$, the
action for the $d$ compact coordinates is
\begin{equation}\label{eq:narain:action}
 S=\frac1{4\pi}\int\!\dd\tau\int_0^{2\pi}\!\dd\sigma\,
 \Big[G_{ij}\big(\dot X^i\dot X^j-X'^iX'^j\big)+2B_{ij}\,\dot X^iX'^j\Big],
\end{equation}
where $B_{ij}=-B_{ji}$ is the two-form in units of $\ap$
\source{GPR §2.4, eq.~(2.4.1), ll.~1031--1053}. The overall sign of the $B$ term is a
convention; the one used here reproduces the canonical momentum of GPR, eq.~(2.4.7), and the
generalized metric fixed in the conventions of this book. For $d=1$ there is no two-form, and
after the rescaling $X\to X/R$ the first term is the circle action of \cref{ch:circle}.

Count the parameters. A symmetric matrix $G$ has $d(d+1)/2$ independent entries and an
antisymmetric matrix $B$ has $d(d-1)/2$, so there are
\begin{equation}\label{eq:narain:count}
 \frac{d(d+1)}2+\frac{d(d-1)}2=d^2
\end{equation}
real parameters \source{GPR §2.4, eqs.~(2.4.3)--(2.4.4), ll.~1068--1089}. They are
conveniently assembled into one matrix, the \emph{background matrix}\index{background matrix}
\begin{equation}\label{eq:narain:E}
 E=G+B,\qquad G=\tfrac12\big(E+E^{\mathsf T}\big),\quad B=\tfrac12\big(E-E^{\mathsf T}\big)
\end{equation}
\source{GPR §2.4, eq.~(2.4.5), ll.~1089--1096}. Each parameter multiplies an operator that
can be added to the worldsheet action without spoiling conformal invariance, because the
action stays quadratic (the model is ``Gaussian''). Such parameters are called
\emph{moduli}\index{moduli}: in the lower-dimensional effective theory each of them is a
massless scalar field without potential. One may wonder whether the term with $B$ matters at
all, since for constant $B$ the integrand $B_{ij}\dot X^iX'^j
=\tfrac12B_{ij}\,\epsilon^{\alpha\beta}\partial_\alpha X^i\partial_\beta X^j$ is a total
derivative. It does not change the equations of motion. It changes the quantum theory of
strings that wind, as we now show.

\section{Canonical quantization of the zero modes}\label{sec:narain:zeromodes}

\subsection{Canonical momentum versus velocity}

From \eqref{eq:narain:action} the momentum density conjugate to $X^i(\sigma)$ is
\begin{equation}\label{eq:narain:P}
 P_i=\frac{\partial\mathcal L}{\partial\dot X^i}
 =\frac1{2\pi}\Big(G_{ij}\dot X^j+B_{ij}X'^j\Big)
\end{equation}
\source{GPR §2.4, eq.~(2.4.7), ll.~1107--1110}. The equations of motion are the free wave
equation $\ddot X^i=X''^i$, and a closed string on the torus must close only up to the
identification \eqref{eq:narain:period}. The general solution is therefore
\begin{equation}\label{eq:narain:modes}
 X^i(\tau,\sigma)=x^i+m^i\sigma+v^i\tau+\text{oscillators},\qquad
 X^i(\tau,\sigma+2\pi)=X^i(\tau,\sigma)+2\pi m^i .
\end{equation}
The integers $m^i$ are the \emph{winding numbers}\index{winding number}: the string wraps
$m^i$ times around the $i$-th cycle (\cref{fig:narain:torus}). The constant $v^i$ is the
centre-of-mass velocity. The total momentum is the integral of \eqref{eq:narain:P}; the
oscillators are periodic and drop out, so
\begin{equation}\label{eq:narain:pzero}
 p_i\equiv\int_0^{2\pi}\!\dd\sigma\,P_i=G_{ij}v^j+B_{ij}m^j .
\end{equation}
Now comes the step that carries all the physics. The operator $p_i$ is conjugate to the
centre-of-mass position $x^i$, $[x^i,p_j]=\ii\,\delta^i_j$, so it generates translations, and a
wavefunction $\ee^{\ii p_ix^i}$ must be single valued under $x^i\to x^i+2\pi$. Hence
\begin{equation}\label{eq:narain:nquant}
 p_i=n_i\in\Z
\end{equation}
\source{GPR §2.4, eq.~(2.4.8), ll.~1110--1115}. It is the \emph{canonical} momentum that is
quantized, not the velocity. Solving \eqref{eq:narain:pzero} for the velocity gives
\begin{equation}\label{eq:narain:velocity}
 v^i=G^{ij}\big(n_j-B_{jk}m^k\big),
\end{equation}
where $G^{ij}$ is the inverse matrix of $G_{ij}$ \source{GPR §2.4, eq.~(2.4.30), ll.~1425--1438}.
A string with winding but $n=0$ therefore still moves when $B\neq0$.

\begin{physicsbox}[Physical picture: $B$ is a theta angle]
A charged particle on a ring threaded by a magnetic flux has canonical momentum
$p=m\dot x+qA$. The vector potential is locally pure gauge and does not enter the classical
motion, but $p$ is quantized in integers, so the allowed \emph{velocities} are shifted by the
flux, and the energy levels depend on it. A flux of one quantum can be removed by relabelling
the integer. The constant $B$-field does the same thing to a winding string: the winding
number $m^k$ plays the role of the charge, $B_{jk}$ the role of the vector potential, and
\eqref{eq:narain:velocity} is the shifted velocity. GPR compare with the theta term of
four-dimensional gauge theory, in whose presence a monopole acquires electric charge
\source{GPR §2.4, ll.~1199--1206}. Two consequences follow at once: $B$ does
matter quantum mechanically, and a shift of $B_{ij}$ by an antisymmetric \emph{integer}
matrix can be absorbed into $n_j\to n_j+\Theta_{jk}m^k$ and changes nothing
\source{GPR §2.4, ll.~1351--1371}. The second fact is the first element of the duality group of
\cref{ch:odd}.
\end{physicsbox}

\subsection{The Hamiltonian}

The Hamiltonian density is $\mathcal H=P_i\dot X^i-\mathcal L$. Inserting
\eqref{eq:narain:P}, the term with $B$ cancels between $P_i\dot X^i$ and $\mathcal L$, which is
the Hamiltonian form of the statement that a total derivative carries no energy:
\begin{equation}\label{eq:narain:Hvel}
 \mathcal H=\frac1{4\pi}\Big(\dot X^{\mathsf T}G\dot X+X'^{\mathsf T}GX'\Big).
\end{equation}
A Hamiltonian must be written in terms of momenta. With $\dot X=G^{-1}(2\pi P-BX')$ from
\eqref{eq:narain:P} and $B^{\mathsf T}=-B$,
\begin{align}
 \dot X^{\mathsf T}G\dot X
 &=(2\pi P-BX')^{\mathsf T}G^{-1}(2\pi P-BX')\nonumber\\
 &=(2\pi)^2P^{\mathsf T}G^{-1}P-X'^{\mathsf T}BG^{-1}BX'+4\pi\,X'^{\mathsf T}BG^{-1}P .
\end{align}
In the second line we used $(BX')^{\mathsf T}=-X'^{\mathsf T}B$ twice: once to turn
$(BX')^{\mathsf T}G^{-1}BX'$ into $-X'^{\mathsf T}BG^{-1}BX'$, and once to combine the two cross
terms, which are equal because each is a number and equals its own transpose. Therefore
\begin{equation}\label{eq:narain:Hdensity}
 \mathcal H=\frac1{4\pi}\Big[(2\pi)^2P^{\mathsf T}G^{-1}P
 +X'^{\mathsf T}\big(G-BG^{-1}B\big)X'+4\pi\,X'^{\mathsf T}BG^{-1}P\Big]
\end{equation}
\source{GPR §2.4, eq.~(2.4.9), ll.~1116--1140}. Notice that $G-BG^{-1}B=G+B^{\mathsf T}G^{-1}B$ is
positive definite, so the energy is bounded below, as it must be.

For the zero modes, $P_i=n_i/2\pi$ and $X'^i=m^i$ are constant along the string; integrating
\eqref{eq:narain:Hdensity} over $\sigma$ multiplies it by $2\pi$ and gives
\begin{equation}\label{eq:narain:Hzero}
 H_0=\frac12\Big[n^{\mathsf T}G^{-1}n+m^{\mathsf T}\big(G-BG^{-1}B\big)m
 +2\,m^{\mathsf T}BG^{-1}n\Big]
\end{equation}
\source{GPR §2.4, eq.~(2.4.11), ll.~1183--1191}. This quadratic form in the $2d$ integers is
best written with a single $2d\times2d$ matrix. Collect winding and momentum numbers in the
\emph{doubled charge vector}\index{doubled charge vector} $Z=(m^i,n_i)$, windings first, and
define
\begin{equation}\label{eq:narain:genmetric}
 \HH(G,B)=\begin{pmatrix}G-BG^{-1}B&BG^{-1}\\[2pt]-G^{-1}B&G^{-1}\end{pmatrix},
 \qquad H_0=\frac12\,Z^{\mathsf T}\HH\,Z
\end{equation}
\source{GPR §2.4, eqs.~(2.4.17)--(2.4.18), ll.~1267--1290, where the matrix is called $M$}.
The two off-diagonal blocks contribute $m^{\mathsf T}BG^{-1}n-n^{\mathsf T}G^{-1}Bm$, and the second
term is the transpose of the first with $B^{\mathsf T}=-B$, so both equal $m^{\mathsf T}BG^{-1}n$,
in agreement with \eqref{eq:narain:Hzero}. The matrix $\HH$ is the \emph{generalized
metric}\index{generalized metric}; it is symmetric and positive definite, and
\cref{ch:genmetric} is devoted to it. The Lean definition \lean{genMetric} (in the file
\texttt{DFT/GeneralizedMetric.lean} of \texttt{DualScaleStream2}) is this matrix, with the
same ordering $Z=(w,n)$.

\section{Left- and right-moving momenta}\label{sec:narain:pLpR}

\subsection{An orthonormal frame}

On the circle the zero modes were packaged into $p_L$ and $p_R$, the momenta of the
left-moving and right-moving parts of $X$. The same is possible on $T^d$, with one new
ingredient: to take the square of a momentum we need an orthonormal frame. Choose a real
invertible $d\times d$ matrix $e$, a \emph{vielbein}\index{vielbein}, with
\begin{equation}\label{eq:narain:vielbein}
 G=e^{\mathsf T}e,\qquad\text{so that}\qquad G^{-1}=e^{-1}e^{-\mathsf T},
\end{equation}
where $e^{-\mathsf T}$ denotes the inverse transpose. Geometrically the columns of
$\sqrt{\ap}\,e$ are the Cartesian components of the lattice vectors $\hat e_i$ of
\eqref{eq:narain:torusdef}. The choice of $e$ is not unique: $e$ and $Oe$ with $O\in O(d)$
give the same $G$. A covector with coordinate components $k_i$ has frame components
$e^{-\mathsf T}k$, and its squared length is $k^{\mathsf T}G^{-1}k$.

\begin{definition}[Left and right momenta]\label{def:narain:pLpR}
For a state with momentum numbers $n\in\Z^d$ and winding numbers $m\in\Z^d$ on the torus with
background matrix $E=G+B$,
\begin{equation}\label{eq:narain:pLpR}
 p_L=\frac1{\sqrt2}\,e^{-\mathsf T}\big(n+E^{\mathsf T}m\big),\qquad
 p_R=\frac1{\sqrt2}\,e^{-\mathsf T}\big(n-E\,m\big),
\end{equation}
both vectors in $\R^d$ with the Euclidean scalar product. Explicitly
$E^{\mathsf T}=G-B$, so $p_L$ contains $n+(G-B)m$ and $p_R$ contains $n-(G+B)m$.
\end{definition}

This is GPR's eq.~(2.4.12) \source{GPR §2.4, ll.~1192--1198} rewritten with column vectors;
GPR normalise their frame by $\sum_ae^a_ie^a_j=2G_{ij}$, which hides the factor
$1/\sqrt2$ that we keep explicit. For $d=1$, $e=R/\sqrt{\ap}$ and
\begin{equation}\label{eq:narain:pLpRcircle}
 p_{L,R}=\frac1{\sqrt2}\Big(\frac{\sqrt{\ap}}R\,n\pm\frac R{\sqrt{\ap}}\,m\Big)
 =\sqrt{\frac\ap2}\,\Big(\frac nR\pm\frac{mR}\ap\Big)
\end{equation}
\source{GPR §2.2, eq.~(2.2.5), ll.~620--649}. The bracket on the right is the dimensionful
$p_{L,R}$ of \cref{ch:circle}; the dimensionless momenta of this chapter are those multiplied
by $\sqrt{\ap/2}$, which are exactly the oscillator zero modes $\tilde\alpha_0$, $\alpha_0$
of that chapter.

Where do \eqref{eq:narain:pLpR} come from? Split the zero-mode part of \eqref{eq:narain:modes}
into a function of $\tau+\sigma$ and a function of $\tau-\sigma$:
$m\sigma+v\tau=\frac12(v+m)(\tau+\sigma)+\frac12(v-m)(\tau-\sigma)$. With
\eqref{eq:narain:velocity},
\[
 G(v+m)=n-Bm+Gm=n+E^{\mathsf T}m,\qquad G(v-m)=n-Bm-Gm=n-Em .
\]
So $p_L$ and $p_R$ are, up to normalisation and the change to the orthonormal frame, the
left-moving and right-moving velocities with the index lowered. This is why they, and not $n$
and $m$ separately, appear in vertex operators and in the Virasoro generators.

\subsection{The two quadratic identities}

\begin{proposition}[Sum and difference of squares]\label{thm:narain:identities}
For all $n,m\in\R^d$, every symmetric invertible $G$ and every antisymmetric $B$,
\begin{align}
 p_L^2+p_R^2&=Z^{\mathsf T}\HH(G,B)\,Z ,\label{eq:narain:sum}\\
 p_L^2-p_R^2&=2\,n\cdot m=Z^{\mathsf T}\eta\,Z ,
 \qquad \eta=\begin{pmatrix}0&\mathbf 1_d\\\mathbf 1_d&0\end{pmatrix},\label{eq:narain:diff}
\end{align}
where $Z=(m,n)$ and $n\cdot m=\sum_in_im^i$. In particular the difference does not depend on
$G$ or $B$.
\end{proposition}

\begin{proof}
By \eqref{eq:narain:vielbein}, $p_L^2=\frac12(n+E^{\mathsf T}m)^{\mathsf T}G^{-1}(n+E^{\mathsf T}m)$
and similarly for $p_R^2$ with $E^{\mathsf T}\to-E$. Expand both:
\begin{align*}
 2p_L^2&=n^{\mathsf T}G^{-1}n+2\,n^{\mathsf T}G^{-1}E^{\mathsf T}m+m^{\mathsf T}EG^{-1}E^{\mathsf T}m,\\
 2p_R^2&=n^{\mathsf T}G^{-1}n-2\,n^{\mathsf T}G^{-1}E\,m+m^{\mathsf T}E^{\mathsf T}G^{-1}E\,m .
\end{align*}
The two quadratic terms in $m$ are equal. Indeed
$EG^{-1}E^{\mathsf T}=(G+B)G^{-1}(G-B)=G-B+B-BG^{-1}B=G-BG^{-1}B$, and
$E^{\mathsf T}G^{-1}E=(G-B)G^{-1}(G+B)$ gives the same result. In the difference they cancel and
the cross terms add up to $2n^{\mathsf T}G^{-1}(E^{\mathsf T}+E)m=4\,n^{\mathsf T}G^{-1}Gm=4\,n\cdot m$;
dividing by $2$ gives \eqref{eq:narain:diff}. In the sum the cross terms give
$2n^{\mathsf T}G^{-1}(E^{\mathsf T}-E)m=-4\,n^{\mathsf T}G^{-1}Bm=4\,m^{\mathsf T}BG^{-1}n$, so
$p_L^2+p_R^2=n^{\mathsf T}G^{-1}n+m^{\mathsf T}(G-BG^{-1}B)m+2m^{\mathsf T}BG^{-1}n$, which is
$2H_0$ of \eqref{eq:narain:Hzero}.
\end{proof}

Both identities, the relation $\HH\eta\HH=\eta$ used below, and $\det\HH=1$ were also
confirmed for $d=2$ with symbolic $G$ and $B$ by computer algebra (a symbolic check, not a
Lean proof). The content of the proposition is that the $2d$ integers carry two quadratic
forms. One, $\HH$, is positive definite and depends on the background: it is the energy. The
other, $\eta$, is indefinite and knows nothing about the background: it will turn out to be
the spin of the state.

\subsection{Mass formula and level matching on \texorpdfstring{$T^d$}{Td}}
\label{sec:narain:massformula}

The oscillators are not affected by the compactification except through their inner product:
in the coordinates $X^i$ the modes obey
$[\alpha^i_k,\alpha^j_l]=[\tilde\alpha^i_k,\tilde\alpha^j_l]=k\,G^{ij}\delta_{k+l,0}$, the
number operators are $N=\sum_{k>0}\alpha^i_{-k}G_{ij}\alpha^j_k$ and the same with tildes,
and the complete Hamiltonian of the compact sector is
\begin{equation}\label{eq:narain:Hfull}
 H=\tfrac12Z^{\mathsf T}\HH Z+N+\tilde N
\end{equation}
\source{GPR §2.4, eqs.~(2.4.32)--(2.4.34), ll.~1448--1489}. For the bosonic string on
$\R^{1,25-d}\times T^d$ the two mass-shell conditions of \cref{ch:circle} hold unchanged, with
$N$ and $\tilde N$ now counting the oscillators of all $24$ transverse directions. In the
dimensionless momenta of this chapter they read
\begin{equation}\label{eq:narain:massLR}
 \frac\ap2M^2=p_R^2+2(N-1),\qquad\frac\ap2M^2=p_L^2+2(\tilde N-1).
\end{equation}
(Multiply the circle formulas $M^2=p_R^2+\frac4\ap(N-1)$ by $\ap/2$ and use
\eqref{eq:narain:pLpRcircle}.) Half the sum and the difference of the two conditions give,
with \cref{thm:narain:identities},
\begin{align}
 M^2&=\frac1\ap\,Z^{\mathsf T}\HH(G,B)Z+\frac2\ap\big(N+\tilde N-2\big),\label{eq:narain:mass}\\
 N-\tilde N&=\tfrac12\big(p_L^2-p_R^2\big)=n\cdot m .\label{eq:narain:level}
\end{align}
These are the $T^d$ versions of the mass formula and of the level-matching
condition\index{level matching} $N-\tilde N=nw$ of \cref{ch:circle}. Level matching is the
statement that the state is invariant under rigid shifts of $\sigma$; it is a constraint on the
integers alone, with no $G$ or $B$ in it. That is the first hint that $n\cdot m$ is more
fundamental than the energy.

It is a common mistake to read \eqref{eq:narain:mass} as ``momentum $n$, energy
$n^{\mathsf T}G^{-1}n$'' and to conclude that $B$ is invisible. That is true only for $m=0$.
Completing the square in \eqref{eq:narain:Hzero},
\begin{equation}\label{eq:narain:square}
 H_0=\tfrac12(n-Bm)^{\mathsf T}G^{-1}(n-Bm)+\tfrac12m^{\mathsf T}Gm ,
\end{equation}
the kinetic energy of the shifted momentum plus the stretching energy of the winding string:
as $B$ varies, a state with fixed labels $(n,m)$ changes its energy continuously.

\section{The Narain lattice}\label{sec:narain:lattice}

\subsection{Lattices with an indefinite inner product}

We need a few definitions from lattice theory; \cref{ch:lattices} has the general theory.

\begin{definition}\label{def:narain:lattice}
A \emph{lattice}\index{lattice} $\Gamma$ of rank $r$ is the set of integer combinations
$\sum_ak^af_a$ of $r$ linearly independent vectors $f_a$ in a real vector space $V$ of
dimension $r$ equipped with a nondegenerate symmetric bilinear form $\langle\ ,\,\rangle$.
Its \emph{Gram matrix}\index{Gram matrix} in the basis $f_a$ is
$\mathcal G_{ab}=\langle f_a,f_b\rangle$. The lattice is \emph{integral} if all
$\langle v,v'\rangle$, $v,v'\in\Gamma$, are integers, and \emph{even}\index{even lattice} if
all $\langle v,v\rangle$ are even integers. The \emph{dual lattice}\index{dual lattice} is
\[
 \Gamma^*=\{u\in V:\ \langle u,v\rangle\in\Z\ \text{for all }v\in\Gamma\},
\]
and $\Gamma$ is \emph{self-dual}\index{self-dual lattice}, or
\emph{unimodular}\index{unimodular lattice}, if $\Gamma^*=\Gamma$. The \emph{signature}
$(n_+,n_-)$ counts the positive and negative eigenvalues of the form on $V$.
\end{definition}

An even lattice is integral, because
$\langle v,v'\rangle=\frac12\big(\langle v+v',v+v'\rangle-\langle v,v\rangle-\langle v',v'\rangle\big)$.
An integral lattice is contained in its dual. Self-duality is decided by one number:

\begin{lemma}\label{thm:narain:det}
An integral lattice is self-dual if and only if its Gram matrix has determinant $\pm1$.
\end{lemma}

\begin{proof}
Let $f^{*a}$ be the basis of $V$ dual to $f_a$, $\langle f^{*a},f_b\rangle=\delta^a_b$. A
vector $u=\sum u_af^{*a}$ pairs integrally with all of $\Gamma$ exactly when all $u_a$ are
integers, so $\Gamma^*$ is the lattice generated by the $f^{*a}$. Expanding
$f_a=\sum_b\mathcal G_{ab}f^{*b}$ shows that $\Gamma\subseteq\Gamma^*$ is the image of $\Z^r$
under the integer matrix $\mathcal G$. This image is all of $\Z^r$ if and only if
$\mathcal G^{-1}$ is again an integer matrix, that is if and only if $\det\mathcal G=\pm1$
(if $\mathcal G^{-1}$ is integral then $\det\mathcal G\cdot\det\mathcal G^{-1}=1$ in $\Z$;
conversely Cramer's rule gives $\mathcal G^{-1}=\operatorname{adj}\mathcal G/\det\mathcal G$).
\end{proof}

The argument in brackets, ``an integer inverse forces $\det=\pm1$'', is precisely the
certificate lemma \lean{isUnimodular\_of\_mul\_eq\_one} of \cref{sec:narain:lean}.

\subsection{Definition and first properties}

Put the left and right momenta of a state side by side, $(p_L,p_R)\in\R^d\oplus\R^d$, and give
this $2d$-dimensional space the \emph{Lorentzian} inner product
\begin{equation}\label{eq:narain:lorentz}
 \big\langle(p_L,p_R),(p_L',p_R')\big\rangle=p_L\cdot p_L'-p_R\cdot p_R' ,
\end{equation}
which has signature $(d,d)$; we write $\R^{d,d}$ for the space with this form.

\begin{definition}[Narain lattice]\label{def:narain:narain}
The \emph{Narain lattice}\index{Narain lattice} of the background $E=G+B$ is
\[
 \Gamma^{d,d}(E)=\big\{\big(p_L(n,m),p_R(n,m)\big):\ n,m\in\Z^d\big\}\subset\R^{d,d},
\]
with $p_{L,R}$ given by \eqref{eq:narain:pLpR}.
\end{definition}

\begin{theorem}[\TL\ in general; \TA\ for the statements listed in
\cref{sec:narain:lean}]\label{thm:narain:evenselfdual}
For every background, $\Gamma^{d,d}(E)$ is an even self-dual lattice of signature $(d,d)$.
In the basis obtained by setting one of the $2d$ integers $(m,n)$ to one and the others to
zero, its Gram matrix is $\eta$ of \eqref{eq:narain:diff}, for every $G$ and $B$.
\end{theorem}

\begin{proof}
The map $(m,n)\mapsto(p_L,p_R)$ is linear and injective (from $p_L-p_R=\sqrt2\,e^{-\mathsf T}Gm$
and then $p_L+p_R$ one recovers $m$ and $n$), so the images of the $2d$ unit vectors are a
basis of a lattice of rank $2d$. Polarizing \eqref{eq:narain:diff},
\begin{equation}\label{eq:narain:pairing}
 \big\langle(p_L,p_R),(p_L',p_R')\big\rangle=n\cdot m'+n'\cdot m=Z^{\mathsf T}\eta Z',
\end{equation}
so the Gram matrix is $\eta$. The norm of a lattice vector is $2\,n\cdot m\in2\Z$: the lattice
is even. Finally $\eta^2=\mathbf 1$, so $\eta$ is its own integer inverse and
$\det\eta=(-1)^d$; by \cref{thm:narain:det} the lattice is self-dual. The signature is that
of $\R^{d,d}$.
\end{proof}

The statement is GPR's \source{GPR §2.4, eq.~(2.4.16), ll.~1249--1258}, where it is
attributed to Narain. For $d=1$ the Gram matrix is
\begin{equation}\label{eq:narain:U}
 U=\begin{pmatrix}0&1\\1&0\end{pmatrix},
\end{equation}
the \emph{hyperbolic plane}\index{hyperbolic plane $U$}: two null vectors, pure momentum
$(n,m)=(1,0)$ and pure winding $(0,1)$, with mutual pairing one. For general $d$ the pairing
\eqref{eq:narain:pairing} only couples $n_i$ with $m^i$ for the same $i$, so
\begin{equation}\label{eq:narain:Usum}
 \Gamma^{d,d}\;\cong\;\underbrace{U\oplus U\oplus\dots\oplus U}_{d},
\end{equation}
one hyperbolic plane per circle. This is no accident of our construction. Even self-dual
lattices of signature $(n_+,n_-)$ exist only for $n_+-n_-$ divisible by $8$, and when both
$n_+$ and $n_-$ are positive the lattice is unique up to isometry
\source{Asp §2.3, ll.~559--566}. For signature $(d,d)$ it is \eqref{eq:narain:Usum}; for
signature $(3,19)$ it is the K3 lattice $3U\oplus2E_8(-1)$ of \cref{ch:k3lattice}, and the
three copies of $U$ in it are the same matrix \eqref{eq:narain:U}.

\begin{pitfallbox}[Common pitfall: ``the lattice depends on the radius'']
As an abstract lattice, a free abelian group with an integer bilinear form, $\Gamma^{d,d}$ is
the same for all $G$ and $B$: it is $\Z^{2d}$ with the form $\eta$. What depends on the
background is the \emph{embedding} in $\R^{d,d}=\R^d_L\oplus\R^d_R$, that is, how each lattice
vector is split into a left and a right part. Both statements are usually abbreviated as
``the Narain lattice'', and one must know which is meant. The spectrum sees the embedding; the
duality group of \cref{ch:odd} is the symmetry group of the abstract lattice.
\end{pitfallbox}

\subsection{The circle: changing \texorpdfstring{$R$}{R} is a Lorentz boost}

For $d=1$ the geometry can be drawn. From \eqref{eq:narain:pLpRcircle}, with
$r=R/\sqrt{\ap}$,
\begin{equation}\label{eq:narain:lightcone}
 p_L+p_R=\sqrt2\,\frac nr,\qquad p_L-p_R=\sqrt2\,m\,r .
\end{equation}
The combinations $p_L\pm p_R$ are light-cone coordinates of the plane $\R^{1,1}$, and changing
the radius rescales one by $r^{-1}$ and the other by $r$. That is a Lorentz boost of rapidity
$\lambda=\ln r$:
\begin{equation}\label{eq:narain:boost}
 \begin{pmatrix}p_L\\p_R\end{pmatrix}_{\!r}
 =\begin{pmatrix}\cosh\lambda&-\sinh\lambda\\-\sinh\lambda&\cosh\lambda\end{pmatrix}
 \begin{pmatrix}p_L\\p_R\end{pmatrix}_{\!r=1},\qquad\lambda=\ln r .
\end{equation}
(Check: at $r=1$, $p_{L,R}=(n\pm m)/\sqrt2$, and the first row gives
$\frac1{\sqrt2}[(\cosh\lambda-\sinh\lambda)n+(\cosh\lambda+\sinh\lambda)m]
=\frac1{\sqrt2}(n/r+mr)$.) A boost preserves $p_L^2-p_R^2$: each lattice point slides along
its hyperbola $p_L^2-p_R^2=2nm$ (\cref{fig:narain:boost}). It does not preserve
$p_L^2+p_R^2$, which is why the masses change. The moduli space of the circle is the group of
boosts, $O(1,1;\R)$ modulo reflections, a half-line parametrised by $r>0$. T-duality is
$\lambda\to-\lambda$: the lattice at $1/r$ is the mirror image under $p_R\to-p_R$ of the
lattice at $r$, with the labels $n\leftrightarrow m$ exchanged.

\begin{figure}[t]
\centering
\begin{tikzpicture}[scale=0.8,>=Stealth]
 \draw[->] (-3.6,0)--(3.8,0) node[right]{$p_L$};
 \draw[->] (0,-3.1)--(0,3.3) node[above]{$p_R$};
 \draw[black!35,dashed] (-3,-3)--(3,3) (-3,3)--(3,-3);
 % hyperbola p_L^2 - p_R^2 = 2 (n m = 1)
 \draw[blue!60!black,domain=-1.35:1.35,samples=60,variable=\t]
   plot ({1.4142*cosh(\t)},{1.4142*sinh(\t)});
 \draw[blue!60!black,domain=-1.35:1.35,samples=60,variable=\t]
   plot ({-1.4142*cosh(\t)},{1.4142*sinh(\t)});
 % hyperbola p_L^2 - p_R^2 = -2 (n m = -1)
 \draw[red!60!black,domain=-1.35:1.35,samples=60,variable=\t]
   plot ({1.4142*sinh(\t)},{1.4142*cosh(\t)});
 \draw[red!60!black,domain=-1.35:1.35,samples=60,variable=\t]
   plot ({1.4142*sinh(\t)},{-1.4142*cosh(\t)});
 % lattice at r = 1 (filled) and r = 1.5 (open)
 \foreach \n in {-2,-1,0,1,2}{\foreach \m in {-2,-1,0,1,2}{
   \pgfmathsetmacro{\pl}{0.7071*(\n+\m)}\pgfmathsetmacro{\pr}{0.7071*(\n-\m)}
   \pgfmathsetmacro{\ql}{0.7071*(\n/1.5+\m*1.5)}\pgfmathsetmacro{\qr}{0.7071*(\n/1.5-\m*1.5)}
   \pgfmathparse{abs(\pl)<3.3 && abs(\pr)<3.0 ? 1 : 0}
   \ifnum\pgfmathresult=1 \fill (\pl,\pr) circle (1.5pt);\fi
   \pgfmathparse{abs(\ql)<3.3 && abs(\qr)<3.0 ? 1 : 0}
   \ifnum\pgfmathresult=1 \draw[thick] (\ql,\qr) circle (2.2pt);\fi}}
 \node[blue!60!black,right] at (2.55,1.75) {\small$nm=1$};
 \node[red!60!black] at (1.35,3.0) {\small$nm=-1$};
 \node[right] at (3.2,-2.4) {\small$\bullet\ r=1$};
 \node[right] at (3.2,-2.9) {\small$\circ\ r=1.5$};
\end{tikzpicture}
\caption{The Narain lattice $\Gamma^{1,1}$ in the $(p_L,p_R)$ plane at the self-dual radius
$r=1$ (filled dots) and at $r=1.5$ (open circles). The dashed lines are the light cone
$p_L^2=p_R^2$, on which the pure-momentum and pure-winding states lie. Changing $r$ is a
Lorentz boost: every point moves along its hyperbola $p_L^2-p_R^2=2nm$. At $r=1$ four points,
$(n,m)=\pm(1,1)$ with $p_R=0$ and $\pm(1,-1)$ with $p_L=0$, sit at distance $\sqrt2$ on the
axes; these are the extra massless states of the self-dual radius (\cref{ch:tduality}).}
\label{fig:narain:boost}
\end{figure}

\section{Why modular invariance forces an even self-dual lattice}
\label{sec:narain:modular}

So far evenness and self-duality were \emph{observed}. The deeper statement is that they are
\emph{required}: a set of momenta $(p_L,p_R)$ defines a consistent closed-string theory only
if it is an even self-dual lattice \source{Asp §3.5, ll.~1985--1988}. The derivation below is
standard; its steps are carried out on the page and have not been checked line by line
against a source in this book's library. The starting point, the one-loop partition function
as a trace over the oscillators and a sum over $(p_L,p_R)$, is
\source{GPR §2.2, eq.~(2.2.13), ll.~766--785}.

\subsection{The torus partition function}

The one-loop vacuum amplitude of a closed string is a path integral over worldsheets with the
topology of a torus (not the target torus $T^d$). A flat worldsheet torus is specified, up to
rescaling, by one complex number $\tau=\tau_1+\ii\tau_2$, $\tau_2>0$: let a closed string of
length $2\pi$ propagate for Euclidean time $2\pi\tau_2$, rotate it by an angle $2\pi\tau_1$, and
glue the ends. Propagation is generated by $H=L_0+\tilde L_0-\frac{c}{12}$ and the rotation by
$L_0-\tilde L_0$, so that, with $q=\ee^{2\pi\ii\tau}$,
\begin{equation}\label{eq:narain:Zdef}
 Z(\tau,\bar\tau)=\Tr\,q^{\tilde L_0-c/24}\,\bar q^{\,L_0-c/24}.
\end{equation}
Here $\tilde L_0=\frac12p_L^2+\tilde N$ and $L_0=\frac12p_R^2+N$ for the $d$ compact bosons,
and $c=d$ (\cref{ch:cft}). The trace factorises into oscillators and zero modes. Each
oscillator $\tilde\alpha_{-k}$ contributes a geometric series
$1+q^k+q^{2k}+\dots=(1-q^k)^{-1}$, so the $d$ left-moving towers give
$q^{-d/24}\prod_{k\ge1}(1-q^k)^{-d}=\eta(\tau)^{-d}$, where
$\eta(\tau)=q^{1/24}\prod_{k\ge1}(1-q^k)$ is the Dedekind eta function\index{Dedekind eta
function} (no relation to the matrix $\eta$). Therefore
\begin{equation}\label{eq:narain:Z}
 Z_\Gamma(\tau,\bar\tau)=\frac{\Theta_\Gamma(\tau,\bar\tau)}{|\eta(\tau)|^{2d}},\qquad
 \Theta_\Gamma(\tau,\bar\tau)=\sum_{(p_L,p_R)\in\Gamma}q^{\frac12p_L^2}\,\bar q^{\frac12p_R^2}
 =\sum_{\Gamma}\ee^{\ii\pi\tau p_L^2-\ii\pi\bar\tau p_R^2}.
\end{equation}
Different values of $\tau$ can describe the same torus. The parallelogram spanned by $1$ and
$\tau$ generates the same lattice in $\C$ as the one spanned by $1$ and $\tau+1$, and, after a
rotation and rescaling, as the one spanned by $1$ and $-1/\tau$. These two maps,
\begin{equation}
 T:\ \tau\mapsto\tau+1,\qquad S:\ \tau\mapsto-\frac1\tau ,
\end{equation}
generate the \emph{modular group}\index{modular group} $\SLZ$ (acting modulo $\pm\mathbf 1$).
The path integral does not know which parallelogram we used, so a consistent theory must have
$Z(\tau+1)=Z(\tau)$ and $Z(-1/\tau)=Z(\tau)$. This is
\emph{modular invariance}\index{modular invariance}. If it failed, the integral over $\tau$
that defines the one-loop amplitude would depend on an arbitrary choice of fundamental
region.

\subsection{\texorpdfstring{$T$}{T} requires an even lattice}

Under $\tau\to\tau+1$, $|\eta(\tau)|$ is unchanged, since $\eta(\tau+1)=\ee^{\ii\pi/12}\eta(\tau)$
is a pure phase. Each term of $\Theta_\Gamma$ is multiplied by
\begin{equation}
 \ee^{\ii\pi(p_L^2-p_R^2)} .
\end{equation}
Terms with different $(p_L^2,p_R^2)$ are linearly independent functions of $\tau$, so the sum
is invariant if and only if every phase is $1$:
\begin{equation}\label{eq:narain:Tcondition}
 p_L^2-p_R^2\in2\Z\quad\text{for all }(p_L,p_R)\in\Gamma .
\end{equation}
This is evenness with respect to the Lorentzian form \eqref{eq:narain:lorentz}. Physically,
$\tau\to\tau+1$ is a rotation of the closed string by a full turn, generated by
$L_0-\tilde L_0$, the worldsheet spin. A full turn must act trivially, so the spin
$N-\tilde N+\frac12(p_R^2-p_L^2)$ of every state, physical or not, must be an integer. For the
toroidal lattice $\frac12(p_L^2-p_R^2)=n\cdot m$: evenness holds because momentum and winding
numbers are integers.

\subsection{\texorpdfstring{$S$}{S} requires a self-dual lattice}

For $S$ we need the Poisson summation formula\index{Poisson resummation}: for a lattice
$\Gamma\subset V$ and a rapidly decreasing function $f$,
\begin{equation}\label{eq:narain:poisson}
 \sum_{v\in\Gamma}f(v)=\frac1{\operatorname{vol}(\Gamma)}\sum_{u\in\Gamma^*}\hat f(u),\qquad
 \hat f(u)=\int_V\dd^{r}x\,f(x)\,\ee^{-2\pi\ii\langle u,x\rangle},
\end{equation}
where $\operatorname{vol}(\Gamma)=\sqrt{|\det\mathcal G|}$ is the volume of a unit cell. (The
formula is usually stated with a Euclidean pairing; the Lorentzian one differs by
$x_R\to-x_R$, which changes neither the measure nor the Gaussian below.) Take
$f(x)=\ee^{\ii\pi\tau x_L^2-\ii\pi\bar\tau x_R^2}$. The Fourier transform is a product of $2d$
one-dimensional Gaussian integrals
$\int\dd x\,\ee^{\ii\pi\tau x^2-2\pi\ii ux}=(-\ii\tau)^{-1/2}\,\ee^{-\ii\pi u^2/\tau}$, and the
complex conjugate formula for the right-movers, so
\[
 \hat f(u)=(-\ii\tau)^{-d/2}(\ii\bar\tau)^{-d/2}\,
 \ee^{\ii\pi(-1/\tau)u_L^2-\ii\pi(-1/\bar\tau)u_R^2},\qquad
 (-\ii\tau)(\ii\bar\tau)=|\tau|^2 .
\]
Inserting this in \eqref{eq:narain:poisson} and renaming $\tau\to-1/\tau$ gives
\begin{equation}\label{eq:narain:thetaS}
 \Theta_\Gamma(-1/\tau)=\frac{|\tau|^{d}}{\operatorname{vol}(\Gamma)}\,\Theta_{\Gamma^*}(\tau).
\end{equation}
The eta function obeys $\eta(-1/\tau)=(-\ii\tau)^{1/2}\eta(\tau)$, hence
$|\eta(-1/\tau)|^{2d}=|\tau|^{d}|\eta(\tau)|^{2d}$, and the factors of $|\tau|$ cancel:
\begin{equation}\label{eq:narain:ZS}
 Z_\Gamma(-1/\tau)=\frac1{\operatorname{vol}(\Gamma)}\,Z_{\Gamma^*}(\tau).
\end{equation}
If $\Gamma=\Gamma^*$ then $\operatorname{vol}(\Gamma)=1$ by \cref{thm:narain:det} and $Z$ is
invariant. Conversely, assume invariance under both $T$ and $S$. By
\eqref{eq:narain:Tcondition} the lattice is even, hence integral, hence
$\Gamma\subseteq\Gamma^*$. Compare the constant terms, $q^0\bar q^0$, of
$\Theta_\Gamma$ and $\Theta_{\Gamma^*}/\operatorname{vol}(\Gamma)$. Only the zero vector has
$p_L=p_R=0$, so the constant terms are $1$ and $1/\operatorname{vol}(\Gamma)$. Thus
$\operatorname{vol}(\Gamma)=1$, the Gram determinant is $\pm1$, and $\Gamma=\Gamma^*$.

\begin{theorem}[\TL]\label{thm:narain:modular}
Let $\Gamma\subset\R^{d,d}$ be a lattice of rank $2d$. The partition function
\eqref{eq:narain:Z} is invariant under $\SLZ$ if and only if $\Gamma$ is even and self-dual.
\end{theorem}

\begin{physicsbox}[Physical picture: what the two conditions mean]
Evenness says that every operator of the theory has integer spin, so that correlation
functions are single valued on the worldsheet. Self-duality says that the set of charges is
as large as it can be: the dual lattice is the set of all charges whose vertex operators are
mutually local with those of $\Gamma$, and $\Gamma^*=\Gamma$ means that none of them has
been left out. A proper sublattice of $\Gamma^{d,d}$, for instance ``momentum states only'',
passes the $T$ test and fails the $S$ test. A point particle on a torus has only momentum
states. Modular invariance, a purely stringy requirement, forces the winding states to be
there. This is the sharpest form of the statement of \cref{ch:circle} that a closed string
cannot be compactified without winding sectors.
\end{physicsbox}

\section{The moduli space}\label{sec:narain:moduli}

\subsection{The generalized metric is an element of \texorpdfstring{$O(d,d;\R)$}{O(d,d;R)}}

The group $\OddR{d}$\index{O(d,d;R)@$O(d,d;\R)$} consists of the real $2d\times2d$ matrices
preserving the form $\eta$,
\begin{equation}\label{eq:narain:odd}
 g^{\mathsf T}\eta\,g=\eta,\qquad g=\begin{pmatrix}a&b\\c&d\end{pmatrix}
 \ \Longleftrightarrow\
 a^{\mathsf T}c+c^{\mathsf T}a=0,\quad b^{\mathsf T}d+d^{\mathsf T}b=0,\quad a^{\mathsf T}d+c^{\mathsf T}b=\mathbf 1
\end{equation}
\source{GPR §2.4, eqs.~(2.4.13)--(2.4.15), ll.~1211--1242}. (In the block matrix, $d$ is a
$d\times d$ block; the clash of letters is traditional.) Its dimension is found by counting:
$g$ has $4d^2$ entries, and $g^{\mathsf T}\eta g=\eta$ is an equation between symmetric matrices,
hence $d(2d+1)$ conditions, leaving
\begin{equation}
 \dim\OddR{d}=4d^2-d(2d+1)=d(2d-1).
\end{equation}

The generalized metric \eqref{eq:narain:genmetric} is itself such a matrix. The quickest way
to see it is to factorise. With the vielbein of \eqref{eq:narain:vielbein} define
\begin{equation}\label{eq:narain:gE}
 g_E=\begin{pmatrix}\mathbf 1&B\\0&\mathbf 1\end{pmatrix}
 \begin{pmatrix}e^{\mathsf T}&0\\0&e^{-1}\end{pmatrix}
 =\begin{pmatrix}e^{\mathsf T}&Be^{-1}\\0&e^{-1}\end{pmatrix}.
\end{equation}
Both factors satisfy \eqref{eq:narain:odd}: for the first, $a=d=\mathbf 1$, $c=0$ and
$b^{\mathsf T}d+d^{\mathsf T}b=B^{\mathsf T}+B=0$; for the second, $b=c=0$ and
$a^{\mathsf T}d=e\,e^{-1}=\mathbf 1$. Multiplying out,
\begin{equation}\label{eq:narain:HgE}
 g_Eg_E^{\mathsf T}
 =\begin{pmatrix}e^{\mathsf T}e+Be^{-1}e^{-\mathsf T}B^{\mathsf T}&Be^{-1}e^{-\mathsf T}\\
 e^{-1}e^{-\mathsf T}B^{\mathsf T}&e^{-1}e^{-\mathsf T}\end{pmatrix}
 =\begin{pmatrix}G-BG^{-1}B&BG^{-1}\\-G^{-1}B&G^{-1}\end{pmatrix}=\HH(G,B)
\end{equation}
\source{GPR §2.4, eqs.~(2.4.20)--(2.4.23), ll.~1294--1325; GPR write $G=ee^{\mathsf T}$, which
exchanges $e$ and $e^{\mathsf T}$}. Since $\OddR{d}$ is closed under transposition
\source{GPR §2.4, ll.~1243--1248}, $\HH=g_Eg_E^{\mathsf T}$ belongs to it:
\begin{equation}\label{eq:narain:HetaH}
 \HH\,\eta\,\HH=\eta,\qquad\text{equivalently}\qquad(\eta\HH)^2=\mathbf 1 .
\end{equation}
The factorisation \eqref{eq:narain:gE} has a plain reading: every toroidal background is
obtained from the square torus at the self-dual radius with no $B$-field, $E=\mathbf 1$,
by a linear change of frame $e\in GL(d,\R)$ followed by a shift of $B$. The counting agrees:
$d^2$ parameters in $e$, plus $d(d-1)/2$ in $B$, minus $d(d-1)/2$ for the rotations
$e\to Oe$ that do not change $G$, is $d^2$.

\subsection{Planes in \texorpdfstring{$\R^{d,d}$}{R(d,d)} and the coset}

Because of \eqref{eq:narain:HetaH} the matrices
\begin{equation}\label{eq:narain:projectors}
 P_\pm=\tfrac12\big(\mathbf 1\pm\eta\HH\big)
\end{equation}
are complementary projectors, $P_\pm^2=P_\pm$, $P_+P_-=0$, $P_++P_-=\mathbf 1$.
\Cref{thm:narain:identities} says what they project on:
\begin{equation}\label{eq:narain:pLsq}
 p_L^2=\tfrac12Z^{\mathsf T}(\HH+\eta)Z=Z^{\mathsf T}\eta P_+Z,\qquad
 p_R^2=\tfrac12Z^{\mathsf T}(\HH-\eta)Z=-Z^{\mathsf T}\eta P_-Z .
\end{equation}
So $P_+$ keeps the left-moving part of a charge vector and $P_-$ the right-moving part. The
image $\Pi$ of $P_+$ is a $d$-dimensional subspace of $\R^{d,d}$ on which the form $\eta$ is
positive definite, the image of $P_-$ is its orthogonal complement $\Pi^\perp$, on which
$\eta$ is negative definite, and $\HH=\eta(P_+-P_-)$ is recovered from $\Pi$. We arrive at
the cleanest description of a toroidal compactification
\source{Asp §3.5, ll.~1989--2006}:

\begin{quote}
\emph{Fix the lattice $\Gamma^{d,d}=\Z^{2d}$ with the form $\eta$ once and for all. A
background is a choice of a positive definite $d$-plane $\Pi\subset\R^{d,d}$. The left and
right momenta of a charge vector are its orthogonal projections on $\Pi$ and $\Pi^\perp$.}
\end{quote}

The group $\OddR{d}$ acts on the set of positive $d$-planes, and the action is transitive
(any orthonormal basis of a positive plane, completed by an orthonormal basis of its
complement, is mapped to any other such basis by an element of the group). The subgroup that
maps a given $\Pi$ to itself consists of the independent rotations of $\Pi$ and of
$\Pi^\perp$, that is, of $p_L$ and of $p_R$, which change neither $p_L^2$ nor $p_R^2$. Hence

\begin{theorem}[\TL]\label{thm:narain:moduli}
The moduli space of toroidal backgrounds is locally
\begin{equation}\label{eq:narain:coset}
 \mathcal M_d=\frac{\OddR{d}}{O(d;\R)\times O(d;\R)},\qquad
 \dim\mathcal M_d=d(2d-1)-2\cdot\frac{d(d-1)}2=d^2 ,
\end{equation}
and $E=G+B\mapsto\HH(G,B)$ is a parametrisation of it.
\end{theorem}
\noindent\source{GPR §2.4, ll.~1207--1210 and 1259--1266}. The dimension agrees with
\eqref{eq:narain:count}, as it must. A third way to get $d^2$: a plane near $\Pi$ is the
graph of a linear map $\Pi\to\Pi^\perp$, a $d\times d$ matrix. This is how the moduli appear
in the spectrum. The states $\alpha^i_{-1}\tilde\alpha^j_{-1}\ket{0}$ with both
oscillators along the torus have $N=\tilde N=1$, $n=m=0$, hence $M^2=0$ by
\eqref{eq:narain:mass}; they are $d^2$ massless scalars in the uncompactified dimensions, and
their vertex operators are the $d^2$ marginal operators
$\partial X^i\bar\partial X^j$ that change $E_{ij}$. For $d=1$ the coset is the half-line of
\cref{fig:narain:boost}; for $d=2$ it has four dimensions and factorises into two copies of
the upper half-plane, the subject of \cref{ch:torus2}.

\begin{pitfallbox}[Common pitfall: ``locally'']
Two different planes can give the same set of pairs $(p_L^2,p_R^2)$, and then the same
physics. If $g$ is an \emph{integer} matrix in $\OddR{d}$, it maps the lattice $\Z^{2d}$ to
itself, and the backgrounds $\HH$ and $g^{\mathsf T}\HH g$ differ only by the relabelling
$Z\to gZ$ of the charges: $Z^{\mathsf T}(g^{\mathsf T}\HH g)Z=(gZ)^{\mathsf T}\HH(gZ)$. The true moduli
space is the quotient $\Odd{d}\backslash\OddR{d}/(O(d)\times O(d))$
\source{Asp §3.5, eq.~(77), ll.~1983--1986}. The relabelling statement, with the bijection
made explicit, is the Lean theorem \lean{spectrum\_equivalence} quoted below; the group
$\Odd{d}$ is the subject of \cref{ch:odd}. GPR write the action as $\HH\to g\HH g^{\mathsf T}$
\source{GPR §2.4, eq.~(2.4.19), ll.~1288--1293}; the two conventions differ by
$g\leftrightarrow g^{\mathsf T}$, which is harmless because the group is closed under
transposition.
\end{pitfallbox}

\begin{historybox}
GPR attribute the even self-dual momentum lattice to Narain (their ref.~[236], 1986), the
coset \eqref{eq:narain:coset} to Narain and to Narain, Sarmadi and Witten (ref.~[238], 1987),
and the matrix \eqref{eq:narain:genmetric} to Giveon, Rabinovici and Veneziano (ref.~[146],
1989) \source{GPR §2.4, ll.~1207--1210, 1249--1258, 1267--1274; refs.\ ll.~10712,
10822--10824}.
\end{historybox}

\section{Worked examples}\label{sec:narain:examples}

\begin{example}[The $SU(3)\times SU(3)$ point of $T^2$]\label{ex:narain:su3}
Take $d=2$ and the background matrix
\begin{equation}\label{eq:narain:su3E}
 E=\begin{pmatrix}1&1\\0&1\end{pmatrix},\qquad
 G=\begin{pmatrix}1&\frac12\\[1pt]\frac12&1\end{pmatrix},\qquad
 B=\begin{pmatrix}0&\frac12\\[1pt]-\frac12&0\end{pmatrix}
\end{equation}
\source{GPR §2.4.2, eq.~(2.4.64), ll.~2026--2036}. Geometrically, the two lattice vectors
have length $\sqrt{\ap}$ and make an angle of $60^\circ$ ($\cos\theta=G_{12}/\sqrt{G_{11}G_{22}}
=\frac12$), and half a unit of $B$-flux threads the torus. We compute
\[
 G^{-1}=\frac13\begin{pmatrix}4&-2\\-2&4\end{pmatrix},\qquad
 BG^{-1}=\frac13\begin{pmatrix}-1&2\\-2&1\end{pmatrix},\qquad
 G-BG^{-1}B=\frac13\begin{pmatrix}4&2\\2&4\end{pmatrix},
\]
so that
\begin{equation}
 \HH=\frac13\begin{pmatrix}4&2&-1&2\\2&4&-2&1\\-1&-2&4&-2\\2&1&-2&4\end{pmatrix},
 \qquad\det\HH=1 .
\end{equation}
Consider the state $(n_1,n_2;m^1,m^2)=(1,0;1,0)$. Then $E^{\mathsf T}m=(1,1)$, $Em=(1,0)$, so
$n+E^{\mathsf T}m=(2,1)$ and $n-Em=(0,0)$. Hence $p_R=0$ and
\[
 p_L^2=\tfrac12\,(2,1)\,G^{-1}\begin{pmatrix}2\\1\end{pmatrix}
 =\tfrac12\cdot\tfrac13\,(16-8+4)=2 ,
\]
consistent with $p_L^2-p_R^2=2\,n\cdot m=2$. By \eqref{eq:narain:level} the state needs
$N-\tilde N=1$; with $N=1$, $\tilde N=0$ the mass formula \eqref{eq:narain:massLR} gives
$\frac\ap2M^2=p_R^2+2(N-1)=0$. If the oscillator $\alpha^\mu_{-1}$ points along the
uncompactified directions this is a \emph{massless vector} carrying winding and momentum.

Enumerating all integers with $|n_i|,|m^i|\le3$ (a finite computer search; larger integers
have larger $p_L^2+p_R^2$ because $\HH$ is positive definite) gives \cref{tab:narain:su3}.
There are exactly six lattice vectors with $(p_L^2,p_R^2)=(2,0)$. In the orthonormal frame
$e=\big(\begin{smallmatrix}1&1/2\\0&\sqrt3/2\end{smallmatrix}\big)$ their left momenta are
\[
 p_L=\pm\sqrt2\,(1,0),\qquad\pm\sqrt2\,\big(\tfrac12,\tfrac{\sqrt3}2\big),\qquad
 \pm\sqrt2\,\big(-\tfrac12,\tfrac{\sqrt3}2\big),
\]
the vertices of a regular hexagon of radius $\sqrt2$: the root system of $SU(3)$. Together
with the two vectors $\alpha^\mu_{-1}\tilde\alpha^i_{-1}\ket0$ that exist at every point of
the moduli space, they fill out the $6+2=8$ gauge bosons of $SU(3)$, and the six vectors with
$(p_L^2,p_R^2)=(0,2)$ give a second $SU(3)$. This generalises the $SU(2)\times SU(2)$ of the
self-dual radius (\cref{ch:tduality}). For comparison, at $E=\mathbf 1$ the same search finds
four vectors of each kind, $\pm(1,0;1,0)$, $\pm(0,1;0,1)$ and $\pm(1,0;-1,0)$,
$\pm(0,1;0,-1)$: the gauge group is $SU(2)^2\times SU(2)^2$
\source{GPR §2.4.2, ll.~2019--2026}. That the extra vectors are non-abelian gauge bosons
is a statement about interactions (\TL); the counting supports it but does not prove it.
\end{example}

\begin{table}[t]
\centering
\caption{The shortest vectors of $\Gamma^{2,2}$ at the point \eqref{eq:narain:su3E}
(exhaustive search over $|n_i|,|m^i|\le3$). Last column: lightest level-matched state, from
$\frac\ap2M^2=\frac12(p_L^2+p_R^2)+N+\tilde N-2$.}
\label{tab:narain:su3}
\begin{tabular}{@{}cccp{0.34\textwidth}c@{}}\toprule
$(p_L^2,p_R^2)$ & number & $n\cdot m$ & examples $(n_1,n_2;m^1,m^2)$ & lightest
$\frac\ap2M^2$\\\midrule
$(0,0)$ & $1$ & $0$ & $(0,0;0,0)$ & $-2$\\
$(\frac23,\frac23)$ & $18$ & $0$ & $(1,0;0,0)$, $(0,0;1,0)$, $(0,1;-1,0)$ & $-\frac43$\\
$(2,0)$ & $6$ & $1$ & $(1,0;1,0)$, $(1,1;0,1)$, $(0,1;-1,1)$ & $0$\ \ $(N=1,\tilde N=0)$\\
$(0,2)$ & $6$ & $-1$ & $(1,0;-1,1)$, $(1,1;-1,0)$, $(0,1;0,-1)$ & $0$\ \ $(N=0,\tilde N=1)$\\
\bottomrule
\end{tabular}
\end{table}


\section{What the machine has checked}\label{sec:narain:lean}

Four Lean files bear on this chapter, written in two independent ``streams'' of the project
that meet here. None of them defines a string, a worldsheet or a partition function: they
contain the arithmetic core of \cref{thm:narain:evenselfdual}, statements about the integers
$(n,w)$ and about explicit integer matrices. All depend only on the axioms \lean{propext},
\lean{Classical.choice} and \lean{Quot.sound}.

\subsection{Stream 1: the rank-two lattice from the momenta}

The file \texttt{StringTheoryFormalization/UseCases/NarainLattice.lean} works over $\Q$ with
$\ap=1$ and uses the unnormalised momenta of \cref{ch:circle} (no factor $1/\sqrt2$), which
it compensates by a factor $\frac12$ in the norm:

\begin{leancode}
def pL (n w R : ℚ) : ℚ := n / R + w * R
def pR (n w R : ℚ) : ℚ := n / R - w * R
def narainNorm (n w R : ℚ) : ℚ := (pL n w R ^ 2 - pR n w R ^ 2) / 2

theorem narain_norm_R_independent (n w R : ℚ) (hR : R ≠ 0) :
    narainNorm n w R = 2 * n * w := by ...
\end{leancode}

\begin{leanbox}[In Lean: the norm forgets the radius]
\textbf{Reading.} For all rationals $n$, $w$ and every nonzero rational $R$,
$\frac12\big[(n/R+wR)^2-(n/R-wR)^2\big]=2nw$. This is \eqref{eq:narain:diff} for $d=1$,
$B=0$, with rational data; here $n$ and $w$ are not even required to be integers.
\textbf{Proof idea.} Unfold the definitions, clear the denominator $R$ with
\texttt{field\_simp} (this is where $R\neq0$ is used), and close the polynomial identity with
\texttt{ring}. \textbf{What it is not.} Nothing is said about $d>1$, $B\neq0$ or irrational radii;
\cref{thm:narain:identities} in that generality is proved on the page, not in Lean.
\end{leanbox}

The file then \emph{forgets} the momenta and keeps the integer quadratic form:

\begin{leancode}
def Q (n w : ℤ) : ℤ := 2 * n * w
theorem narain_form_even (n w : ℤ) : Even (Q n w) := ⟨n * w, by unfold Q; ring⟩

def B (n w n' w' : ℤ) : ℤ := n * w' + n' * w
theorem B_is_polarization_of_Q (n w n' w' : ℤ) :
    2 * B n w n' w' = Q (n + n') (w + w') - Q n w - Q n' w' := by ...

def gram : Matrix (Fin 2) (Fin 2) ℤ :=
  !![B 1 0 1 0, B 1 0 0 1; B 0 1 1 0, B 0 1 0 1]
theorem gram_eq_hyperbolic : gram = !![0, 1; 1, 0] := by ...
theorem narain_gram_unimodular : gram.det = -1 := by ...
\end{leancode}

\begin{leanbox}[In Lean: even, and Gram determinant $-1$]
\textbf{Reading.} (i) For all integers $n,w$ the integer $2nw$ is even; Mathlib's
\lean{Even a} means $\exists r,\ a=r+r$, and the witness supplied is $r=nw$. (ii) The
bilinear form \lean{B} (no relation to the Kalb--Ramond field; the name is local to this
file) is the polarization of \lean{Q}; this is \eqref{eq:narain:pairing} for $d=1$. (iii) The
Gram matrix of \lean{B} in the basis $(1,0),(0,1)$ is the matrix $U$ of \eqref{eq:narain:U},
and the determinant of this one matrix is $-1$. \textbf{Proof idea.} \texttt{ring} for (i)
and (ii); \texttt{norm\_num} evaluates the four entries; \texttt{Matrix.det\_fin\_two} is the
$2\times2$ determinant formula. \textbf{What it is not.} \lean{narainNorm} and \lean{Q} are linked only by the shared
expression $2nw$; no Lean theorem says ``the set of pairs $(p_L,p_R)$ is a lattice isometric
to $U$'', and the step from $\det=-1$ to $\Gamma=\Gamma^*$ is \cref{thm:narain:det}, proved
on the page.
\end{leanbox}

\subsection{Stream 2: Gram matrices, a certificate lemma and the hyperbolic plane}

\texttt{DualScaleStream2/Lattice/Basic.lean} represents a lattice of rank $n$ by its Gram
matrix alone and defines the two properties of this chapter as predicates on it:

\begin{leancode}
abbrev Gram (n : ℕ) := Matrix (Fin n) (Fin n) ℤ
def IsEvenDiag {n : ℕ} (G : Gram n) : Prop := ∀ i, Even (G i i)
def IsUnimodular {n : ℕ} (G : Gram n) : Prop := G.det = 1 ∨ G.det = -1

theorem isUnimodular_of_mul_eq_one {n : ℕ} (G H : Gram n) (h : H * G = 1) :
    IsUnimodular G := by ...

theorem even_quadratic_form_of_even_diag {n : ℕ} (G : Gram n)
    (hsymm : Gᵀ = G) (heven : IsEvenDiag G) (x : Fin n → ℤ) :
    Even (x ⬝ᵥ (G *ᵥ x)) := by ...
\end{leancode}

\begin{leanbox}[In Lean: two general lemmas about integer Gram matrices]
\textbf{Reading.} Both hold for every rank $n$. The first says that an integer matrix with
an integer left inverse has determinant $\pm1$; it is one half of \cref{thm:narain:det}, and
Stream~2 discharges every unimodularity claim this way, by exhibiting an inverse. The second
says that a symmetric integer matrix with even diagonal gives an even number $x^{\mathsf T}Gx$
for \emph{every} integer vector $x$: evenness need only be checked on a basis. \textbf{Proof idea.} For the first, $\det H\cdot\det G=\det\mathbf 1=1$
in $\Z$, and the only ways to write $1$ as a product of two integers are $1\cdot1$ and
$(-1)(-1)$. For the second, split $G=D+T+T^{\mathsf T}$ into its diagonal and strictly upper
triangular parts; then $x^{\mathsf T}Gx=\sum_iG_{ii}x_i^2+2\,x^{\mathsf T}Tx$ and both terms are even.
\textbf{What it is not.} \lean{IsUnimodular} is a statement about a determinant. The
equivalence with $\Gamma^*=\Gamma$ is not formalised, because the dual lattice is not defined
in these files.
\end{leanbox}

\texttt{DualScaleStream2/Lattice/Hyperbolic.lean} then posits the hyperbolic plane as a
literal matrix and checks its properties entry by entry:

\begin{leancode}
def hyperbolicU : Gram 2 := !![0, 1; 1, 0]

theorem hyperbolicU_symm : hyperbolicUᵀ = hyperbolicU := by ...
theorem hyperbolicU_evenDiag : IsEvenDiag hyperbolicU := by ...
theorem hyperbolicU_mul_self : hyperbolicU * hyperbolicU = 1 := by ...
theorem hyperbolicU_unimodular : IsUnimodular hyperbolicU :=
  isUnimodular_of_mul_eq_one _ _ hyperbolicU_mul_self
theorem hyperbolicU_congruence :
    (!![1, 1; 1, -1] : Gram 2)ᵀ * hyperbolicU * !![1, 1; 1, -1]
      = !![2, 0; 0, -2] := by ...

theorem hyperbolicU_eq_narain_gram :
    hyperbolicU = StringTheory.UseCases.NarainLattice.gram := by ...
\end{leancode}

\begin{leanbox}[In Lean: $U$ is even and unimodular, and the two streams agree]
\textbf{Reading.} The matrix $U$ is symmetric, has even diagonal, squares to the identity and
is therefore unimodular by the certificate lemma; with
\lean{even\_quadratic\_form\_of\_even\_diag}, $x^{\mathsf T}Ux$ is even for all $x\in\Z^2$. In the
congruence $P^{\mathsf T}UP=\diag(2,-2)$ the columns of $P$ are the charge vectors $(1,1)$ and
$(1,-1)$ of norm $\pm2$, the special points of \cref{fig:narain:boost}; concluding
``signature $(1,1)$'' from it uses Sylvester's law of inertia, which is \TL\ here. The last
theorem says that the matrix written down in Stream~2 and the matrix \emph{computed} in
Stream~1 by polarizing $Q(n,w)=2nw$ are equal. \textbf{Proof idea.} There are only four entries; \texttt{fin\_cases} enumerates them
and each case closes by \texttt{rfl}, \texttt{simp} or \texttt{norm\_num}. \textbf{Why the bridge matters.} The $U$ that enters the K3 lattice $3U\oplus2E_8(-1)$ in
\cref{ch:k3lattice} is, by a kernel-checked equality, the momentum--winding Gram matrix of
one circle, not merely a matrix that looks the same in two files.
\end{leanbox}

\subsection{Stream 2: general \texorpdfstring{$d$}{d}}

For arbitrary $d$ the relevant declarations are in
\texttt{DualScaleStream2/TDuality/Factorized.lean} and \texttt{Spectrum.lean}. The index type
of a doubled charge vector is \lean{Charge d}, defined as \lean{Fin d ⊕ Fin d}, and
\lean{eta d} is the block matrix \lean{fromBlocks 0 1 1 0}, the $\eta$ of
\eqref{eq:narain:diff}.

\begin{leancode}
def chargeNorm (Z : Charge d → ℤ) : ℤ := Z ⬝ᵥ (eta d *ᵥ Z)

theorem chargeNorm_sumElim (n w : Fin d → ℤ) :
    chargeNorm (Sum.elim n w) = 2 * (n ⬝ᵥ w) := by ...
theorem chargeNorm_even (Z : Charge d → ℤ) : Even (chargeNorm Z) := by ...
theorem chargeNorm_invariant (g : Matrix (Charge d) (Charge d) ℤ)
    (hg : IsODD g) (Z : Charge d → ℤ) :
    chargeNorm (g *ᵥ Z) = chargeNorm Z := by ...

theorem spectrum_equivalence (g : Matrix (Charge d) (Charge d) ℤ)
    (hg : IsODD g) (H : Matrix (Charge d) (Charge d) ℝ) :
    ∃ e : (Charge d → ℤ) ≃ (Charge d → ℤ), ∀ Z : Charge d → ℤ,
      massForm ((toReal g)ᵀ * H * toReal g) (fun i => (Z i : ℝ)) =
        massForm H (fun i => (e Z i : ℝ)) ∧
      chargeNorm (e Z) = chargeNorm Z := by ...
\end{leancode}

\begin{leanbox}[In Lean: $\Gamma^{d,d}$ is even for every $d$; dual backgrounds match charge
by charge]
\textbf{Reading.} \lean{chargeNorm Z} is $Z^{\mathsf T}\eta Z$. For every $d$ and every integer
vector $Z$, split into its two blocks, it equals $2\,n\cdot w$, hence is even: the evenness
half of \cref{thm:narain:evenselfdual} for all $d$. (Since $\eta$ is symmetric it does not
matter which block is called momentum; the docstrings say $Z=(n,w)$, while \lean{massForm}
is used with windings first, as in \eqref{eq:narain:genmetric}.) \lean{IsODD g} is the
equation $g^{\mathsf T}\eta g=\eta$ for an integer matrix. \lean{spectrum\_equivalence} says: for
every such $g$ and \emph{every} real matrix $H$, not only a generalized metric, there is a
bijection $e$ of $\Z^{2d}$, namely $Z\mapsto gZ$, with
$Z^{\mathsf T}(g^{\mathsf T}Hg)Z=(eZ)^{\mathsf T}H(eZ)$ and equal charge norms: the content of the
pitfall box ``locally'' in \cref{sec:narain:moduli}. \textbf{Proof idea.} The inverse of $g$
is the integer matrix $\eta g^{\mathsf T}\eta$ (\lean{isODD\_left\_inv},
\lean{isODD\_right\_inv}), which gives the bijection; the mass identity is associativity of
matrix multiplication after casting $\Z\to\R$ (\lean{massForm\_covariant}); the norm identity
is $g^{\mathsf T}\eta g=\eta$. \textbf{What it is not.} Only the zero-mode quadratic form is
concerned: no oscillators, no partition function. Unimodularity of $\eta$ for general $d$
would follow from \lean{eta\_mul\_self} ($\eta^2=\mathbf 1$) by the certificate argument, but is
not stated as an \lean{IsUnimodular} theorem, because \lean{Gram n} is indexed by \lean{Fin n}
and \lean{eta d} by \lean{Fin d ⊕ Fin d} (\cref{exc:narain:leaneta}).
\end{leanbox}

\subsection{Tier table and what a faithful formalization would need}

\Cref{tab:narain:tiers} collects the status of every claim of this chapter.
\begin{table}[!t]
\centering\small
\caption{Epistemic status of the claims of this chapter.}\label{tab:narain:tiers}
\begin{tabular}{@{}p{0.50\textwidth}cp{0.38\textwidth}@{}}\toprule
Claim & Tier & Evidence\\\midrule
$\frac12(p_L^2-p_R^2)=2nw$ for $d=1$, $B=0$, rational $n,w,R$, unnormalised momenta & \TA &
\lean{narain\_norm\_R\_independent}\\
$2nw$ is even; its polarization; Gram matrix $=U$; $\det U=-1$ & \TA &
\lean{narain\_form\_even}, \lean{B\_is\_polarization\_of\_Q}, \lean{gram\_eq\_hyperbolic},
\lean{narain\_gram\_unimodular}\\
$U$ symmetric, even diagonal, $U^2=\mathbf 1$, unimodular, congruent to $\diag(2,-2)$ & \TA &
\lean{hyperbolicU\_symm}, \lean{\_evenDiag}, \lean{\_mul\_self}, \lean{\_unimodular},
\lean{\_congruence}\\
Stream~2's $U$ equals Stream~1's Gram matrix & \TA & \lean{hyperbolicU\_eq\_narain\_gram}\\
Integer inverse $\Rightarrow\det=\pm1$; even diagonal $\Rightarrow$ even form (any rank) &
\TA & \lean{isUnimodular\_of\_mul\_eq\_one}, \lean{even\_quadratic\_form\_of\_even\_diag}\\
$Z^{\mathsf T}\eta Z=2n\cdot w$ is even and $\Odd{d}$-invariant, all $d$ & \TA &
\lean{chargeNorm\_sumElim}, \lean{chargeNorm\_even}, \lean{chargeNorm\_invariant}\\
$H$ and $g^{\mathsf T}Hg$ have the same zero-mode spectrum, all $d$ & \TA &
\lean{spectrum\_equivalence}\\
\Cref{thm:narain:identities} (general $d$, $G$, $B$); \cref{thm:narain:det} & --- & proved
on the page; the former also checked symbolically for $d=2$\\
Signature of $U$ is $(1,1)$ (Sylvester) & \TL & quoted in \texttt{Hyperbolic.lean}\\
Modular invariance $\Leftrightarrow$ even self-dual (\cref{thm:narain:modular}) & \TL &
Asp §3.5; derivation on the page\\
Moduli space $\OddR{d}/O(d)\times O(d)$; uniqueness of $\Gamma^{d,d}$ & \TL & GPR §2.4; Asp
§2.3\\
Identification of $(\Z^2,U)$ with the charge lattice of a physical string & --- & modelling
step; never Tier~A\\\bottomrule
\end{tabular}
\end{table}

A faithful formalization of this chapter would need, in increasing order of difficulty: (i)
\cref{thm:narain:identities} over $\R$ with a positive definite $G$ and its inverse, within
reach of Mathlib's matrix library; (ii) lattices as $\Z$-submodules of a real quadratic space
with their duals, so that ``self-dual'' means $\Gamma^*=\Gamma$ rather than $\det=\pm1$;
(iii) the theta series \eqref{eq:narain:Z} and Poisson summation in indefinite signature, for
\cref{thm:narain:modular}; (iv) the coset \eqref{eq:narain:coset} as a manifold. Items (iii)
and (iv) are research-level projects (\cref{ch:open}).

\begin{summarybox}
\begin{itemize}[leftmargin=*,itemsep=2pt]
\item A torus $T^d$ with constant $G$ and $B$ has $d^2$ parameters $E=G+B$. The canonical
 momentum is quantized, $p_i=n_i\in\Z$, the velocity is $G^{-1}(n-Bm)$: $B$ acts on winding
 strings like a theta angle, and integer shifts of $B$ are invisible.
\item $p_L=\frac1{\sqrt2}e^{-\mathsf T}(n+E^{\mathsf T}m)$, $p_R=\frac1{\sqrt2}e^{-\mathsf T}(n-Em)$ with
 $G=e^{\mathsf T}e$. Two identities hold: $p_L^2+p_R^2=Z^{\mathsf T}\HH Z$ (energy, depends on the
 background) and $p_L^2-p_R^2=2\,n\cdot m=Z^{\mathsf T}\eta Z$ (spin, does not).
\item $M^2=\frac1\ap Z^{\mathsf T}\HH Z+\frac2\ap(N+\tilde N-2)$ and $N-\tilde N=n\cdot m$.
\item The vectors $(p_L,p_R)$ form the Narain lattice $\Gamma^{d,d}\cong U^{\oplus d}$: even,
 self-dual, signature $(d,d)$, Gram matrix $\eta$ for every background. Only its embedding in
 $\R^d_L\oplus\R^d_R$ varies; for $d=1$ varying $R$ is a Lorentz boost.
\item Modular invariance requires exactly this: $T$ forces evenness (integer spin), $S$ forces
 self-duality (Poisson resummation). A lattice of momentum states alone fails.
\item A background is a positive $d$-plane in $\R^{d,d}$; the moduli space is locally
 $\OddR{d}/O(d)\times O(d)$, of dimension $d^2$, and globally a quotient by $\Odd{d}$.
\item Lean has checked the rank-two lattice in two independent ways and that they agree,
 two general lemmas on integer Gram matrices, and evenness, invariance and the spectrum
 bijection for all $d$; partition functions and the coset are not formalised.
\end{itemize}
\end{summarybox}

\section*{Exercises}
\addcontentsline{toc}{section}{Exercises}

\begin{exercise}[$\star$]
Show that $\dim\OddR{d}=d(2d-1)$ by linearising $g=\mathbf 1+\epsilon X$ in
\eqref{eq:narain:odd}: the condition becomes $(\eta X)^{\mathsf T}=-\eta X$. Deduce
\eqref{eq:narain:coset}, and list $\dim\mathcal M_d$ for $d=1,2,4,6$. Which of these numbers
is the number of massless scalars $\alpha^i_{-1}\tilde\alpha^j_{-1}\ket0$?
\end{exercise}

\begin{exercise}[$\star$]
Derive \eqref{eq:narain:Hvel} from \eqref{eq:narain:action} and \eqref{eq:narain:P}, and then
the completed square \eqref{eq:narain:square}. Use it to show
that $B\to B+\Theta$, $n\to n+\Theta m$ leaves $H_0$ and $n\cdot m$ invariant for every
antisymmetric integer matrix $\Theta$. For $d=2$, $G=\mathbf 1$, $B_{12}=b$, follow the
states $(n;m)=(0,1;1,0)$ and $(0,0;1,0)$ from $b=0$ to $b=-1$ and show that they exchange
their energies $1$ and $\frac12$. Write the $4\times4$ matrix $g$ acting on $Z=(m,n)$ and
check \eqref{eq:narain:odd}.
\end{exercise}

\begin{exercise}[$\star\star$]
Prove $\HH\eta\HH=\eta$ by block multiplication, without using the factorisation
\eqref{eq:narain:gE}. Then show that $P_\pm$ of \eqref{eq:narain:projectors} are projectors
of rank $d$ (hint: $\Tr(\eta\HH)=\Tr(BG^{-1})-\Tr(G^{-1}B)=0$) and that their images are
orthogonal with respect to $\eta$.
\end{exercise}

\begin{exercise}[$\star\star$]
At $E=\mathbf 1$ for $d=2$, find all $(n;m)$ with $(p_L^2,p_R^2)=(2,0)$ and all with $(1,1)$,
and compute the mass of the lightest level-matched state in each class. Compare with
\cref{tab:narain:su3}. Then move to $G=\mathbf 1$, $B_{12}=\frac12$: do any of the $(2,0)$
vectors survive?
\end{exercise}

\begin{exercise}[$\star\star$]
For an integer $k\ge2$ let $\Gamma_k=\{(n,m):\ m\in k\Z\}\subset\Gamma^{1,1}$ (only
windings divisible by $k$). Compute its Gram matrix, $\operatorname{vol}(\Gamma_k)$ and
$\Gamma_k^*$. Which condition of \cref{thm:narain:modular} fails, what is
$Z_{\Gamma_k}(-1/\tau)$ according to \eqref{eq:narain:ZS}, and which vectors of $\Gamma_k^*$
violate evenness?
\end{exercise}

\begin{exercise}[$\star\star$, Lean]
In the namespace of \texttt{NarainLattice.lean} state and prove: (a) \lean{Q n w = Q w n},
the invariance of the norm under $n\leftrightarrow w$; (b) the symmetry
\lean{B n w n' w' = B n' w' n w}; (c) \lean{Q n w = B n w n w}. Each proof is one line:
\texttt{unfold} followed by \texttt{ring}.
\end{exercise}

\begin{exercise}[$\star\star\star$, Lean]\label{exc:narain:leaneta}
Define the Gram matrix of $U\oplus U$ as a literal \lean{Gram 4} in the ordering
$(m^1,m^2,n_1,n_2)$, so that it equals $\eta$ for $d=2$. Prove that it is unimodular using
\lean{isUnimodular\_of\_mul\_eq\_one} with the matrix itself as the inverse, and that
$x^{\mathsf T}Gx$ is even for all $x$ using \lean{even\_quadratic\_form\_of\_even\_diag}. Then
explain what would be needed to state the same result for \lean{eta d} with $d$ arbitrary
(hint: Mathlib's \lean{Matrix.reindex} along \lean{finSumFinEquiv}, an equivalence
\lean{Fin d ⊕ Fin d ≃ Fin (d + d)}).
\end{exercise}

\section*{Further reading}
\addcontentsline{toc}{section}{Further reading}
\begin{itemize}[leftmargin=*,itemsep=2pt]
\item Toroidal backgrounds, $p_{L,R}$, the matrix $M=\HH$ and the coset:
 \source{GPR §2.4, ll.~1031--1345}; the circle, its partition function and the dimensionless
 metric: \source{GPR §2.2, ll.~576--821}; oscillators under duality:
 \source{GPR §2.4, ll.~1423--1520}; enhanced symmetry:
 \source{GPR §2.4.1--2.4.2, ll.~1836--1866, 2019--2036}.
\item The Grassmannian picture and modular invariance: \source{Asp §3.5, ll.~1965--2011};
 even self-dual lattices: \source{Asp §2.3, ll.~541--566}. Toroidal duality among abelian and
 non-abelian dualities, and the caveat that duals of momentum operators are winding operators
 only for flat compact isometries: \source{AAL §1, ll.~105--135; §2.1, ll.~605--612}.
\item In this book: the lattice theory is developed in \cref{ch:lattices}; the discrete group
 $\Odd{d}$ and its generators in \cref{ch:odd}; the case $d=2$ in \cref{ch:torus2}; the
 generalized metric as a field in \cref{ch:genmetric}; the same coset with signature
 $(4,20)$ in \cref{ch:stringsk3}.
\end{itemize}
